\section{Globular models and the weak-equivalence test}
\label{sec:preliminaries}

We first fix the algebraic conventions and prove the elementary homotopy
calculus needed below.  No model structure is used in this section.

\subsection{The fixed coherator and its free models}

A \emph{globular set} consists of sets of cells in dimensions
$0,1,\ldots$, with source and target maps satisfying the globular identities:
the source and target of a positive-dimensional cell have the same lower
source and the same lower target.  Two cells of the same dimension are
\emph{parallel} if their immediate sources and targets agree; any two
$0$-cells are parallel.  For a cell of dimension greater than $i$, write
$s_i,t_i$ for its $i$-dimensional source and target.  The representable
globular set $d_j$ is the $j$-globe, and $\partial d_j$ is its boundary;
$\partial d_0=\varnothing$.

A \emph{globular sum shape}, or \emph{pasting scheme}, is a globe or a
finite iterated amalgamation
\begin{equation}\label{eq:prelim-shape}
 S=d_{h_1}\amalg_{d_{b_1}}d_{h_2}
       \amalg_{d_{b_2}}\cdots\amalg_{d_{b_{l-1}}}d_{h_l},
 \qquad b_a<h_a,h_{a+1},
\end{equation}
where each gluing identifies a target face with a source face.  Its height
$\dim(S)$ is its largest cell dimension.  A tuple of shape $S$ in a
globular set $X$ means a globular map $S\to X$.  Let $\Theta_0$ denote the
category of these shapes and their globular maps.  Its designated sums are
the amalgamations in Equation~\eqref{eq:prelim-shape}.  The language of
globular theories and their pasting arities goes back to
Berger \cite[Definition~1.5]{Berger2002}.

\begin{definition}[The coherator convention]\label{def:coherator}
A \emph{globular theory} is a category under $\Theta_0$, with the same
objects and preserving the designated globular sums.  A
\emph{Grothendieck coherator} $\mathcal C$ is such a theory with the
following two properties.
\begin{enumerate}
\item It admits a defining tower
$\mathcal C_0=\Theta_0\to\mathcal C_1\to\cdots$ with
$\mathcal C\cong\colim_{r\ge0}\mathcal C_r$.  At stage $r+1$, one
freely adjoins a filler for each pair in a chosen set of parallel pairs
$d_{j-1}\rightrightarrows S$ in $\mathcal C_r$ satisfying
$j>0$ and $\dim(S)\le j$.  The new arrow $d_j\to S$ has that pair
as its source and target.
\item Every parallel pair $d_{j-1}\rightrightarrows S$ with
$\dim(S)\le j$ has a filler $d_j\to S$.
\end{enumerate}
Here freeness includes globular completion: the designated sums continue
to be sums after the new arrows have been adjoined.
\end{definition}

This is the admissibility convention of Ara
\cite[Sections~2.3 and 2.9--2.11]{Ara2013}, fixed in
\cite[Section~5.3]{Henry2016}.  The height bound concerns the
\emph{generating} fillers; a composite operation of output dimension $j$
can have an arity of height greater than $j$.  We retain the chosen free
cellular order throughout the proof and never replace it by an order on
output dimensions.  The tower index records the stage at which an
operation is adjoined, independently of its output dimension.  This
sequential presentation of the theory is separate from the possibly
transfinite cellular presentations of its models below.

Fix $\mathcal C$ and put $\mathcal A=\Mod(\mathcal C)$.  Its objects are
presheaves $\mathcal C^{\op}\to\mathbf{Set}$ taking the designated sums
to limits.  Equivalently, a model is a globular set with an interpretation
of each freely adjoined operation on its tuple set, subject to its two
specified face equations.  There are no additional equations equating
independently chosen fillers.  One obtains this description successively
from the universal property of each free stage: an interpretation of a
new arrow is exactly a function on tuples with the prescribed faces, and
globular completion interprets sums as tuple sets.

To make this universal-property argument explicit, let $G$ be a globular
set and let $G(S)$ be its set of tuples of shape $S$.  Form a category
$\mathcal E_G$ with objects labelled by the shapes and
\[
 \Hom_{\mathcal E_G}(S,T)=\Hom_{\mathbf{Set}}(G(T),G(S)).
\]
Composition is composition of functions in the indicated opposite
direction.  Restriction of tuples defines $\Theta_0\to\mathcal E_G$.
The designated globular sums are sums in $\mathcal E_G$: their tuple
sets are the corresponding finite limits, and maps from any fixed set
preserve those limits.  An interpretation at stage $r$ is therefore a
functor $\mathcal C_r\to\mathcal E_G$ under $\Theta_0$.  A function
interpreting a new filler, with its two prescribed faces, extends this
functor by the free-adjunction and globular-completion universal
properties \cite[Section~2.9]{Ara2013}.  Passage up the tower gives
the model.  Equations among composites follow from functoriality;
no independent filler equation is being assumed.

The forgetful functor to globular sets has a left adjoint $F$.  Set
\[
 D_j=F(d_j),\qquad \partial D_j=F(\partial d_j),\qquad
 I_j:\partial D_j\longrightarrow D_j,\qquad
 J_j:D_j\xrightarrow{s}D_{j+1}.
\]
Put $I=\{I_j\mid j\ge0\}$ and $J=\{J_j\mid j\ge0\}$.
Thus maps $D_j\to X$ are $j$-cells, and maps $\partial D_j\to X$ are
specified boundaries.  The representable presheaf $\mathcal C(-,S)$ is
a model, since the designated diagrams are sums in $\mathcal C$.
Yoneda identifies its maps into $X$ with $X(S)$, the set of tuples of
shape $S$.  Consequently
\begin{equation}\label{eq:prelim-free-shape}
 F(S)=\mathcal C(-,S),\qquad
 \Hom_{\mathcal A}(D_j,F(S))=\mathcal C(d_j,S).
\end{equation}
In particular, cells of $F(S)$ are precisely operations on the arity $S$.

An \emph{$I$-cellular extension} is a transfinite composite of pushouts
of coproducts of the $I_j$; an object is \emph{cellular} if its map from
the initial model is such an extension.  ``Finitely cellular'' means
finitely many attachments, including $I_0:\varnothing\to D_0$.\footnote{We
include the initial vertex attachment.  Although Definition~5.3.2 of
\cite{Henry2016} prints $n>0$, Section~5.3.1 specifies the empty boundary
$S_{-1}$, and Proposition~5.3.7 starts with finite globular sets and the
marked object $D_0$.  Our indexing makes that convention explicit.}
For a globular set $K$, adjoining its cells in increasing dimension is
a globular boundary-cell presentation; applying $F$ makes $F(K)$
cellular.  This does not assert that its underlying globular set has
only the cells of $K$.

The algebraic constructions below are instances of finite-limit sketch
theory; see \cite[Section~4.4, Theorems~1, 2 and~4]{BarrWells1985}.
We recall the specific finiteness facts needed here.  Models of this small
finitary essentially algebraic theory form a locally presentable
category, hence have the required limits and colimits.  Concretely,
free constructions use finite terms in the specified operations with
the prescribed matching and face equations.  Filtered colimits are
computed on the sets of cells: finite tuples and their finite matching
equations commute with filtered colimits of sets.  It follows from
adjunction that $F(K)$ is finitely presentable for every finite globular
set $K$, in particular for the disks and their boundaries.  Thus the
small-object argument of Quillen \cite{Quillen1967}, in the cellular
form of \cite[Theorem~2.1.14 and Corollary~2.1.15]{Hovey1999}, is
available for any subfamily of $I$ or $J$: adjoin cells for all lifting
problems, repeat, and take the
sequential colimit.  Each lifting problem into that colimit already
occurs at a finite stage.  These are algebraic constructions and make
no homotopical assertion.

\subsection{Homotopies, division, and change of boundary}

Contractibility of the theory supplies operations of the following
forms.  There are units $e(a):a\to a$; write $e_p^q(a)$ for their
iteration from dimension $p$ to $q$, and set $e_p^p(a)=a$.  For
$r$-cells $\alpha,\gamma$ and $i<r$ with $s_i\gamma=t_i\alpha$,
there are compositions $\gamma*_i\alpha$ and $i$-inverses
$w_i(\alpha)$.  Their exact face equations at level $i$ are
\[
 s_i(\gamma*_i\alpha)=s_i\alpha,\quad
 t_i(\gamma*_i\alpha)=t_i\gamma,\qquad
 s_iw_i(\alpha)=t_i\alpha,\quad t_iw_i(\alpha)=s_i\alpha.
\]
For $i<k<r$, their $k$-faces are obtained by applying the same
operations to the corresponding $k$-faces.  To construct them, start
in dimension $i+1$ with these prescribed faces and then proceed in
increasing dimension, prescribing the already constructed operations
on each face.  The arities are $d_r\amalg_{d_i}d_r$ and $d_r$, of
height $r$, so each required filler is admissible.  Units are fillers
of the pair consisting twice of the same operation.  Fix all these
choices once in $\mathcal C$.

Write $\alpha\sim\beta$ when there is one $(r+1)$-cell from an
$r$-cell $\alpha$ to a parallel $r$-cell $\beta$.  Units, inverses and
top-level compositions make this an equivalence relation, also for
$r=0$.  Put
\[
 \pi_0(X)=X_0/{\sim},\qquad
 \pi_r(X;u,v)=\{\alpha:u\to v\}/{\sim}\quad(r\ge1),
\]
where $u,v$ are parallel $(r-1)$-cells.  Top-level composition makes
these sets the hom sets of a groupoid on the $(r-1)$-cells having a
fixed lower boundary.  Indeed associativity, unit and inverse
expressions are parallel on the appropriate string of globes or on
a single globe; an admissible $(r+1)$-filler makes their classes equal.
Composition at any level $i$ respects homotopy: compose the given
homotopy with the unit on the fixed factor in dimension $r+1$.

The expansion argument will require a cell with an arbitrary prescribed
boundary, whereas based homotopy groups use iterated unit boundaries.
Division and change of boundary connect these two kinds of tests.
The point requiring care is that multiplication below the top level
can alter the entire boundary.  The following is the division
calculus of \cite[Lemma~4.12]{Ara2013}; we give its boundary correction
explicitly.

\begin{lemma}[Division with specified boundaries]\label{lem:prelim-division}
Fix an $r$-cell $\gamma$ and $i<r$.  Left multiplication
$\alpha\mapsto\gamma*_i\alpha$ induces a bijection from the
homotopy set at any prescribed input boundary composable with $\gamma$
to the homotopy set
at its induced output boundary.  The same holds for right
multiplication.  Empty homotopy sets are allowed.
\end{lemma}

\begin{proof}
We describe universal expressions before evaluating them in a model.
For each $k>i$, let $(g,a)$ be the generic pair of $i$-composable
$k$-cells on $d_k\amalg_{d_i}d_k$.  Initially put
\[
 U_{i+1}^k(g,a)=w_i(g)*_i(g*_i a).
\]
We inductively construct $U_h^k$ for $k\ge h>i$ with the following
two properties: its faces below $h$ equal those of $a$, and its
$l$-faces for $h\le l<k$ are $U_h^l$ evaluated on the corresponding
$l$-faces of $(g,a)$.  These statements hold for $h=i+1$ by the
exact face equations just given.  In particular $U_h^h(g,a)$ and
$a$ are parallel \emph{as operations on the generic pair of
$h$-globes}.  Choose admissible operations
\[
 c_h(g,a):a\longrightarrow U_h^h(g,a),\qquad
 d_h(g,a):U_h^h(g,a)\longrightarrow a.
\]
They have output dimension $h+1$ on an arity of height $h$.
For $k>h$, define
\begin{equation}\label{eq:prelim-normalization}
 \begin{split}
 U_{h+1}^k(g,a)={}&
 e_{h+1}^k d_h(t_hg,t_ha)\ *_h\bigl(
 U_h^k(g,a)\ *_h e_{h+1}^k c_h(s_hg,s_ha)\bigr).
 \end{split}
\end{equation}
The two compositions are defined because the old $h$-faces are
$U_h^h$ on the indicated face inputs.  Their new $h$-faces are
exactly $s_ha,t_ha$, and lower faces retain their required values.
Faces above $h$ obey the same formula in their own dimension,
by the composition and unit face equations.  This proves both
inductive properties.  Notice that parallelism for $c_h,d_h$ is
an identity in their arity, and does not result from a special
equality in the model under consideration.

Now prescribe the input boundary of an $r$-cell $\alpha$, without
assuming that such a cell exists, and let $\beta$ have the induced
boundary of $\gamma*_i\alpha$.  Begin with
$z=w_i(\gamma)*_i\beta$.  For $h=i+1,\ldots,r-1$, replace $z$ by
the right side of Equation~\eqref{eq:prelim-normalization}, with
$k=r$, $U_h^r(g,a)$ replaced by $z$, and the correction factors
evaluated at $(s_h\gamma,s_h\alpha)$ and
$(t_h\gamma,t_h\alpha)$.  These faces of $\alpha$ belong to the
\emph{prescribed boundary}.  At every stage the proper faces of
$z$ obey the universal formulas already proved, because
$\beta$ has the prescribed output boundary.  The construction
therefore finishes with an $r$-cell $K_\gamma(\beta)$ having
exactly the desired input boundary.

All correction factors depend only on $\gamma$ and that boundary.
Write $L_\gamma$ for left multiplication by $\gamma$ on homotopy classes.
Thus $K_\gamma$ is a composite of multiplication maps by fixed
cells and induces a map of homotopy sets.  On an actual input
$\beta=\gamma*_i\alpha$, it equals $U_r^r(\gamma,\alpha)$.
The latter is parallel to $\alpha$ on the generic pair of
$r$-globes, so an admissible $(r+1)$-filler shows that
$K_\gamma L_\gamma=\id$ on classes.

The same construction with left and right interchanged gives a
left inverse for every right multiplication.  Consequently all
left and right multiplication maps are injective.  Each factor
of $K_\gamma$ is one of these now-known injective maps, so
$K_\gamma$ is injective as well.  Applying this injectivity to
$K_\gamma L_\gamma K_\gamma=K_\gamma$ gives
$L_\gamma K_\gamma=\id$.  This proves bijectivity without a
nonemptiness assumption on the input set.
\end{proof}

For a vertex $x$, the group
\[
 \pi_r(X,x)=
 \pi_r\bigl(X;e_0^{r-1}(x),e_0^{r-1}(x)\bigr),\qquad r\ge1,
\]
uses top-level composition.  These are the globular homotopy groups
in Henry's statement.  A model map induces group homomorphisms,
since it preserves the chosen operations.

\begin{lemma}[Change of boundary and vertex]\label{lem:prelim-boundary-change}
The endomorphism homotopy set at any $(r-1)$-cell $u$ is naturally
in bijection with the underlying set of $\pi_r(X,x)$, where $x$
is either its source or its target vertex.  An arrow $x\to y$
gives bijections between the underlying sets of $\pi_r(X,x)$
and $\pi_r(X,y)$, natural in model maps.  If
$\pi_r(X;u,v)$ is nonempty, a chosen cell $u\to v$ gives a
bijection from the endomorphism set at $u$ to this homotopy set.
\end{lemma}

\begin{proof}
The last assertion is composition in the groupoid above.  When
$r=1$, the first assertion is the definition and change of
vertex is conjugation in the fundamental groupoid.  For $r\ge2$,
put $x=s_0u$, $a=e_0^{r-1}(x)$ and $w=u*_0a$, and abbreviate
$e(u)$ to $1_u$.  Lemma~\ref{lem:prelim-division} gives a span
of bijections
\[
 \pi_r(X;u,u)
 \xrightarrow{\ [\alpha]\mapsto[\alpha*_0 1_a]\ }
 \pi_r(X;w,w)
 \xleftarrow{\ [\beta]\mapsto[1_u*_0\beta]\ }
 \pi_r(X;a,a).
\]
For the target vertex $y=t_0u$, use
$b=e_0^{r-1}(y)$ and $w'=b*_0u$ instead: the maps are
$[\alpha]\mapsto[1_b*_0\alpha]$ and
$[\beta]\mapsto[\beta*_0 1_u]$.  Finally, for an arrow
$p:x\to y$, apply these two comparisons to
$u=e_1^{r-1}(p)$.  Naturality follows from the fixed choices of
operations; inverses of the resulting bijections are natural too.
Only set bijections are needed here.
\end{proof}

\subsection{The exact-boundary criterion}

Let $\mathcal W$ consist of the maps inducing a bijection on
$\pi_0$ and isomorphisms on all $\pi_r$ at all source vertices.
The next formulation is \cite[Theorem~4.18]{Ara2013}, with the
change-of-boundary part supplied by
\cite[Theorem~4.13 and Corollaries~4.14--4.15]{Ara2013}.

\begin{proposition}[Weak-equivalence test]\label{prop:weak-test}
For a model map $f:X\to Y$, the following are equivalent.
\begin{enumerate}
\item $f\in\mathcal W$.
\item It induces a bijection on $\pi_0$ and a bijection
\[
 \pi_r(X;u,v)\longrightarrow\pi_r(Y;fu,fv)
\]
for every $r\ge1$ and every parallel pair $u,v$ in $X$.
\item For every $r\ge0$, every specified boundary
$b:\partial D_r\to X$, and every $r$-cell $y$ of $Y$ with
boundary $fb$, there are an $r$-cell $x$ with boundary
\emph{exactly} $b$ and an $(r+1)$-cell $fx\to y$ in $Y$.
\end{enumerate}
In particular, a weak equivalence reflects homotopy between
parallel cells, including equality of components for vertices.
\end{proposition}

\begin{proof}
Suppose (1) holds.  Lemma~\ref{lem:prelim-boundary-change} shows
that $f$ is bijective on every endomorphism homotopy set.
We prove (2) by induction on $r$.  First address existence of
a class: if a $1$-cell joins $fu$ to $fv$, injectivity on
$\pi_0$ gives a cell $u\to v$.  If $r>1$ and a target
$r$-cell joins $fu$ to $fv$, then $fu,fv$ represent the same
class in the appropriate $(r-1)$-dimensional homotopy set.
Injectivity already proved in dimension $r-1$ gives a source
$r$-cell $u\to v$.  Thus a nonempty target homotopy set has
a nonempty source homotopy set.  Choose a source joining cell;
composition with it and its image identifies the map on these
homotopy sets with the map on endomorphism sets.  The latter
is bijective.  If the target set is empty, the source is empty
as well, so the assertion also holds in that case.

Condition (2) immediately gives (3), using surjectivity on
$\pi_0$ when $r=0$.  Conversely (3) gives surjectivity on
$\pi_0$ and on every homotopy set at a source boundary.
If parallel source $r$-cells have homotopic images, apply
(3) in dimension $r+1$ to their exact boundary and the target
joining cell.  This produces a source joining cell.  For
$r=0$ it proves injectivity on $\pi_0$, and for $r\ge1$
it proves injectivity on the indicated homotopy sets.  Hence
(2) follows, and its based instances give (1).
\end{proof}

\begin{lemma}[Closure properties]\label{lem:weak-closure}
The class $\mathcal W$ has two-out-of-three, is closed under
retracts, and is closed under filtered colimits of arrows.
Consequently a transfinite composite of weak equivalences is
a weak equivalence.
\end{lemma}

\begin{proof}
For two-out-of-three, use the induced maps of components and
based groups.  In the case where $f$ and $gf$ are weak
equivalences, a vertex of the intermediate model is joined to
some $f(x)$, by surjectivity of $\pi_0(f)$.
Lemma~\ref{lem:prelim-boundary-change} transfers bijectivity
of the homomorphism induced by $g$ at $f(x)$ to that vertex.
The other two cases follow directly by cancellation of
bijections.  Retract closure follows from
Proposition~\ref{prop:weak-test}(3), even for a retract of
underlying globular arrows: apply its disk and homotopy
witnesses in the larger arrow and retract them.

For a filtered colimit of weak equivalences, the prescribed
boundary, target disk, and their finitely many matching
equations in Proposition~\ref{prop:weak-test}(3) all occur
together at some stage.  Apply the test there and map its
two witnesses to the colimit.  Finally, transfinite induction
handles a transfinite composite: use composition at successor
stages and, at a limit, apply filtered-colimit closure to the
arrows from the constant initial stage to the preceding stages.
No injectivity of the transition maps is required.
\end{proof}

\begin{definition}\label{def:prelim-contractible}
A model $X$ is \emph{contractible} if $X\to1$ belongs to
$\mathcal W$.  Equivalently, $X$ is nonempty and every
parallel pair of cells has a joining cell.
\end{definition}

The equivalence in this definition is
Proposition~\ref{prop:weak-test}(3) applied to $X\to1$:
the condition in dimension zero is nonemptiness, and in
positive dimension it is exact boundary filling.  Equivalently,
$X\to1$ lifts against every $I_j$.  Thus a map from any
cellular domain extends across an $I$-cellular extension
whenever its target is contractible.

For a set of maps $K$, write $K\text{-}\mathrm{inj}$ for the maps
with the right lifting property against every member of $K$.

\begin{lemma}[Lifting tests]\label{lem:lifting-test}
If $f\in\mathcal W$ has the right lifting property against
$J_j$, it has the right lifting property against $I_j$.
Moreover
\[
 I\text{-}\mathrm{inj}
   =\mathcal W\cap J\text{-}\mathrm{inj}.
\]
Every model maps to the terminal model with the right lifting
property against all $J_j$.
\end{lemma}

\begin{proof}
For a square against $I_j$, Proposition~\ref{prop:weak-test}
provides an approximate $j$-cell lift $x$ and a target
$(j+1)$-cell $H:fx\to y$.  Lifting $H$ against
$J_j:D_j\to D_{j+1}$ with prescribed source $x$ gives a
$(j+1)$-cell whose target is a strict lift of $y$ with the
specified boundary.  This also proves the assertion for
$j=0$.

An $I$-injective map satisfies the test strictly, using the
unit on the target disk as its homotopy, so it belongs to
$\mathcal W$.  Each $J_j$ is the composite of two boundary
attachments: first adjoin the parallel target $j$-cell using
$I_j$, then the joining $(j+1)$-cell using $I_{j+1}$.
Thus $I$-injectivity implies $J$-injectivity, and the first
paragraph proves the reverse inclusion in the displayed
equality.  Finally a $J_j$-lifting problem for $X\to1$
is solved by the unit on its prescribed source cell.
\end{proof}

The same unit defines a retraction $D_{j+1}\to D_j$ of
$J_j$.  Its pushout along any specified cell therefore has
a retraction as well.  A transfinite composite of such
pushouts has compatible retractions to its initial object:
choose the retraction at each successor and use the induced
map out of the colimit at each limit.  The existence of these
retractions alone does not establish weak equivalence.
